% geometry/part_intro.tex

% \part{Geometry and Curvature}

\chapter*{Introduction}
\addcontentsline{toc}{chapter}{Introduction}

\begin{remark}
Railway alignment geometry is traditionally described through curves,
radii, transition elements, and engineering rules.
\end{remark}

\begin{remark}
Within AIM, geometry is interpreted differently:
not primarily as collections of geometric primitives, but as controlled
curvature evolution along arc-length.
\end{remark}

\begin{remark}
This shifts the focus:
\begin{itemize}
\item from coordinates to intrinsic structure,
\item from stored geometry to reconstruction,
\item and from isolated curve elements to coupled state evolution.
\end{itemize}
\end{remark}

\begin{remark}
The central geometric quantity is therefore not position itself, but
curvature:
\[
\kappa(s).
\]
\end{remark}

\begin{remark}
Within the AIM framework:
\begin{itemize}
\item geometry emerges from curvature integration,
\item transition curves become curvature-transition laws,
\item and sparse alignment models become compact curvature programs.
\end{itemize}
\end{remark}

\begin{remark}
This geometric viewpoint naturally connects:
\begin{itemize}
\item reconstruction,
\item realization,
\item runtime evaluation,
\item dynamics,
\item and optimization.
\end{itemize}
\end{remark}

\begin{remark}
The chapters of this part introduce the intrinsic geometric structure of
railway alignments, beginning with curvature space and continuing toward
transition classification and higher-order geometric interpretation.
\end{remark}

\begin{proposition}
Within AIM, railway geometry is fundamentally interpreted as curvature
evolution along intrinsic arc-length.
\end{proposition}
